\documentclass[11pt,a4paper]{article}
\usepackage[margin=22mm,top=19mm,bottom=20mm]{geometry}
\usepackage[T1]{fontenc}
\usepackage{lmodern,microtype}
\usepackage[hidelinks,pdftitle={Rule42}]{hyperref}
\pagestyle{plain}
\setlength{\parindent}{0pt}
\setlength{\parskip}{7pt}
\emergencystretch=1em

\begin{document}
{\LARGE\bfseries Rule42}\par
{\small Working design / 1 November 2026}

\textbf{Shared research question}\par
\emph{Can evolutionary search learn compact adapters that recover safe task
completion around frozen world-model systems, reducing reliance on manually
tuned behavioural rules?}

\bigskip
{\Large\bfseries Experiment 1: learning safe behaviour around LeWM}\par
{\small Page 1: method and environments}

\textbf{Purpose.} Learn a small behavioural adapter while keeping the pretrained
world model and task planner frozen. The aim is useful safe completion with less
task-specific tuning. Matching an expert-designed reference could be valuable;
beating it is not required. Explicit safety definitions are still needed.

\textbf{Three environments.} TwoRoom tests navigation and provides the first
mechanism analysis. PushT tests moving and orienting the T-block. OGBench Cube
tests moving the cube to its original object goal, not merely reaching with the
gripper. LeWM is the opening study before the main Block Plucking study;
ManiGuard is deferred.

\textbf{What should carry across.} Reuse the adapter family and learning recipe,
training new weights for each environment. Test fresh cases within each task.
Declare differences in units, action dimensions and task goals. This does not
claim that the same weights transfer between different pretrained models.

\textbf{Task and hazards.} Keep each original task goal. Add supplied forbidden
regions and no-hazard controls. Specify exactly what is protected: the navigation
body, pusher, gripper or other named geometry. Cube needs an explicit vertical
hazard extent. A safe gripper path alone does not prove a feasible object-moving
sequence. Admit cases by a fixed procedure, not by favourable controller outcomes.

\textbf{Adapter.} Qualify two possible interfaces on development cases: choose an
intermediate visual target, or make a bounded correction to the planner's command.
The residual may be simpler to reuse across tasks and Block Plucking. Select one
family before the main comparison. Do not yet fix the earlier TwoRoom-specific
112-parameter actor as the common architecture.

\textbf{Learning and safety.} All learners use the same actor and information
within a task, with generic initialisation and no expert fallback. A passive
monitor records actual violations but never vetoes actions. Remove inherited
hazard masks and shields. A learned predictive safety checker is unnecessary
unless it has a declared role. Training may be unsafe in simulation.

\textbf{Evolutionary method.} Start with CMA-ES for a compact neural actor. Evaluate
task return and episode violation frequency separately. Prefer candidates meeting
a stated empirical risk limit; among them, prefer higher return. Rank failing
candidates by safety shortfall. This avoids a tuned safety-penalty multiplier,
but small samples do not certify the true risk.

\newpage
{\Large\bfseries Experiment 1: comparisons, ablations and the dial}\par
{\small Page 2: evidence and scope}

\textbf{Main comparisons.} Across TwoRoom, PushT and Cube, compare CMA-ES,
parameter random search and PPO-Lagrangian learning the same adapter. Include the
frozen planner, an untrained adapter and an expert-designed fixed reference.
PPO-Lagrangian adjusts its safety multiplier during learning; it still has
hyperparameters. Count every method's development and selection work.

\textbf{Shared objective.} Use the verified task reward and the same risk target.
For RL, charge one safety cost at the first violation and zero thereafter; its
undiscounted episode total records whether any violation occurred. Monitor the
whole executed action prefix, including between adapter decisions. Preserve
original task goals and correct termination for every method.

\textbf{Focused ablations.} Use TwoRoom for the fuller analysis: trained versus
untrained; feasibility-based versus task-only evolution; feasibility-based versus
penalised evolution; training without hazard information; and removing history
or simplifying the actor if these components remain. If using target selection,
examine memory and readout contributions. Retrain removal ablations. Repeat the
essential learning and safety-selection checks in PushT and Cube.

\textbf{Keep the workload focused.} A fixed-penalty PPO sweep can test tuning
sensitivity in TwoRoom. Extra optimizers, every architecture in every task,
shielding and a continuous dial are not all required at once. A shielded
deployment comparison is separate from the main learning result.

\textbf{Safety dial.} First test one fixed risk target. Then train or select a small
library of controllers at different requested violation-probability limits.
Initially test this in TwoRoom. Keep hazards, protected geometry and the violation
definition fixed. Stricter settings request less risk; looser settings permit
risk without requiring it. Lock policies and dial levels before final evaluation.

\textbf{What would establish a useful dial?} Measure actual risk, completion and
engagement at each setting, with uncertainty. Check whether stricter requests
reduce risk while retaining useful performance. Do not assume monotonicity.
Identical policies, uniformly zero violations or growing inactivity alone do not
establish the trade-off. Changing clearance instead is a different specification
dial. One network conditioned on risk can be a later extension.

\textbf{Evaluation.} Use fresh source cases and repeated training runs. Report safe
completion, all completion, violations, unfinished episodes and engagement with
denominators. Separate training risk from final performance. Match interaction
allowances and report actual simulator work, model queries and time. Keep tasks
and physical hosts identifiable; repeated cases are not independent new data.

\textbf{Automation and next steps.} Freeze the recipe after development and count
all remaining tuning. Define acceptable performance differences before claiming
comparability. First qualify task correctness and adapter authority; then freeze
the main comparison, focused ablations and a separately scoped dial study.
Budgets, seeds, sample sizes and stopping rules remain open. This is a working
design, not an execution authorisation.

{\small Full technical detail: \href{FirstExperiment.md}{\texttt{FirstExperiment.md}}.}
\end{document}
